\documentclass[reqno,11pt]{amsart}
\usepackage{math327-lecture}
\usepackage{needspace}
\newcommand{\C}{\mathbb{C}}
\newcommand{\F}{\mathbb{F}}
\newcommand{\Nm}{\operatorname{N}}
\lectureinfo{32}{Dec. 2, 2026}{Irreducibles, Primes, and Unique Factorization Domains}

% Typeset from handwritten Lecture 34, Dec. 3, 2025, pp. 1--11.
\begin{document}
\raggedbottom
\makelecturetitle

\section{Units, associates, primes, and irreducibles}

Last time we met the inclusions
\[
  \{\text{Euclidean domains}\}\subset\{\text{PIDs}\}\subset\{\text{UFDs}\},
\]
where UFD stands for \emph{unique factorization domain}. Today we
define UFDs and prove the second inclusion. We begin with some
preliminaries. Throughout, $R$ is a commutative ring with $1$.

\begin{define}[Units and associates]
An element $u\in R$ is a \emph{unit} if $u$ has an inverse in $R$.
Two elements $a,b\in R$ are \emph{associates}, written $a\sim b$,
if $a=bu$ for some unit $u$.
\end{define}

% NOTE: the handwritten definitions omit the standing conditions
% "nonzero and not a unit" for primes and irreducibles; these are
% the standard conventions and are added here.
\begin{define}[Prime and irreducible elements]
Let $p,x\in R$ be nonzero nonunits.
\begin{itemize}
  \item $p$ is \emph{prime} if whenever $p\mid ab$, either $p\mid a$
  or $p\mid b$.
  \item $x$ is \emph{irreducible} if whenever $x=ab$, either $a$ or
  $b$ is a unit (and then the other factor is an associate of $x$).
\end{itemize}
\end{define}

If $R$ is an integral domain (that is, $ab=0$ implies $a=0$ or
$b=0$), then every prime element is irreducible. The converse holds
in a PID.

\begin{remark}
Divisibility can be phrased in terms of ideals:
$a\mid b$ if and only if $(b)\subseteq(a)$.
\end{remark}

\begin{lemma}\label{lem:irred-prime}
If $R$ is a PID, then every irreducible element of $R$ is prime.
\end{lemma}
\begin{proof}
Let $q$ be irreducible with $q\mid ab$, and assume $q\nmid a$.
We want to show $q\mid b$. Since $R$ is a PID,
\[
  (q,a)=(d)\qquad\text{for some }d\in R.
\]
Because $(q)\subseteq(d)$ and $(a)\subseteq(d)$, we have $d\mid q$
and $d\mid a$. As $q$ is irreducible, $d\mid q$ forces $d$ to be
either a unit or an associate of $q$. If $d\sim q$, then $d\mid a$
gives $q\mid a$, contrary to our assumption. Hence $d$ is a unit,
that is, $(q,a)=(1)$. So
\[
  rq+sa=1\qquad\text{for some }r,s\in R.
\]
Multiplying by $b$ gives
\[
  rqb+sab=b.
\]
Since $q\mid ab$ by assumption, $q$ divides the left-hand side, so
$q$ divides the right-hand side, that is, $q\mid b$.
\end{proof}

\section{Unique factorization domains}

What does it mean to factor an element $x\in R$? If $x$ is not
irreducible, write $x=ab$ with neither factor a unit. If $a$ or $b$
is not irreducible, factor it further, say $a=a'a''$ and $b=b'b''$,
and continue:
\[
  x=ab=(a'a'')(b'b'')=\cdots.
\]
Does this process terminate?

\begin{define}[Unique factorization domain]
An integral domain $R$ is called a \emph{unique factorization
domain} (UFD) if
\begin{enumerate}[label=(\arabic*)]
  \item factorization into irreducibles terminates, so every nonzero
  nonunit is a product of irreducibles; and
  \item whenever
  \[
    x=p_1\cdots p_s=q_1\cdots q_\ell
  \]
  with all $p_i,q_j$ irreducible, we have $s=\ell$ and, after
  reordering, $p_i\sim q_i$ for every $i$.
\end{enumerate}
\end{define}

\begin{example}
The integers $\Z$ form a UFD.
\end{example}

\begin{proposition}\label{prop:pid-ufd}
Every PID is a UFD.
\end{proposition}
\begin{proof}
For now, assume that factorization terminates in PIDs; we return
to this point in Section~\ref{sec:termination}. We prove uniqueness.
Suppose $a\in R$ has two factorizations
\[
  a=p_1\cdots p_r=q_1\cdots q_s
\]
with all $p_i,q_j$ irreducible. Since $R$ is a PID,
Lemma~\ref{lem:irred-prime} shows that the $p_i$ and $q_j$ are also
prime.

The prime $p_1$ divides $q_1\cdots q_s$, so $p_1\mid q_j$ for some
$j$. But $q_j$ is irreducible, so $p_1\sim q_j$; write $q_j=up_1$
with $u$ a unit. Since we are in a domain, we may cancel $p_1$ and
obtain
\[
  p_2\cdots p_r=u\,q_1\cdots\widehat{q_j}\cdots q_s,
\]
where the hat means that $q_j$ is omitted. Absorbing $u$ into one
of the remaining factors, which keeps it irreducible, we get two
factorizations with fewer factors. Repeating (or inducting on $r$)
gives $r=s$ and pairs the factors off into associates.
\end{proof}

\section{Termination of factorization}\label{sec:termination}

\begin{proposition}\label{prop:acc}
Let $R$ be an integral domain. The following are equivalent.
\begin{enumerate}[label=(\arabic*)]
  \item Factorization into irreducibles terminates.
  \item $R$ does not contain a strictly increasing infinite chain of
  principal ideals
  \[
    (a_1)\subsetneq(a_2)\subsetneq(a_3)\subsetneq(a_4)\subsetneq\cdots.
  \]
\end{enumerate}
\end{proposition}
\begin{proof}
We show that (2) implies (1). Let $a_1\in R$, and assume that the
factorization of $a_1$ does not terminate. Then we can write
\[
  a_1=a_2b_2,\qquad a_2=a_3b_3,\qquad a_3=a_4b_4,\qquad\ldots,
\]
where at each stage neither factor is a unit and the factorization
of $a_{n+1}$ again does not terminate. Since $a_{n+1}\mid a_n$, we
have $(a_n)\subseteq(a_{n+1})$; equality would make $b_{n+1}$ a unit
(cancel in the domain), so $(a_n)\subsetneq(a_{n+1})$. This gives
\[
  (a_1)\subsetneq(a_2)\subsetneq(a_3)\subsetneq\cdots,
\]
contradicting (2). The converse is similar and left as an exercise.
\end{proof}

\begin{proposition}\label{prop:pid-terminates}
If $R$ is a PID, then factorization terminates.
\end{proposition}
\begin{proof}
By Proposition~\ref{prop:acc}, it suffices to rule out an infinite
strictly increasing sequence of principal ideals
\[
  (f_1)\subsetneq(f_2)\subsetneq(f_3)\subsetneq\cdots.
\]
% NOTE: the handwritten notes say "sequence of prime ideals" here;
% "principal ideals" is meant.

\emph{Claim:} $J=\bigcup_{n=1}^\infty(f_n)$ is also an ideal.
Indeed, if $a,b\in J$, then since the ideals are nested there is a
single $n$ with $a,b\in(f_n)$; hence $a-b\in(f_n)\subseteq J$.
Similarly $ra\in J$ for all $r\in R$.

Since we are in a PID, $J=(b)$ for some $b$. But $b\in J$, so
$b\in(f_n)$ for some $n$. Then
\[
  (b)\subseteq(f_n)\subsetneq(f_{n+1})\subsetneq\cdots\subseteq J=(b),
\]
which contradicts the strictly increasing nature of the sequence of
ideals.
\end{proof}

This completes the proof of Proposition~\ref{prop:pid-ufd}, so
\[
  \{\text{Euclidean domains}\}\subset\{\text{PIDs}\}\subset\{\text{UFDs}\}.
\]

\begin{example}
If $F$ is a field, then $F[x]$ is a PID, and therefore $F[x]$ is
also a UFD.
\end{example}

\begin{warningbox}
The ring $\Z[x]$ is \emph{not} a PID. It is nevertheless a UFD, as
we will see next time.
\end{warningbox}

\section{A non-example: \texorpdfstring{$\Z[\sqrt{-5}]$}{Z[sqrt(-5)]}}

Consider
\[
  \Z[\sqrt{-5}]=\{a+b\sqrt{-5}:a,b\in\Z\},
\]
with addition and multiplication inherited from
$\Z[\sqrt{-5}]\subset\C$. In this ring,
\[
  6=2\cdot3=(1+\sqrt{-5})(1-\sqrt{-5}).
\]

\emph{Claim:} these are not equivalent factorizations into
irreducibles.

\begin{remark}
Compare with $\Z[i]$, where
\[
  5=(2+i)(2-i)=(1+2i)(1-2i).
\]
These two factorizations \emph{are} equivalent: for example,
$1+2i=i(2-i)$ and $1-2i=-i(2+i)$.
\end{remark}

To show that $2$, $3$, $1+\sqrt{-5}$, and $1-\sqrt{-5}$ are
irreducible and that $2$ and $3$ are not associates of
$1\pm\sqrt{-5}$, introduce the norm.

\begin{definition}[Norm]
For $\alpha=m+n\sqrt{-5}\in\Z[\sqrt{-5}]$, define
\[
  \Nm(\alpha)=m^2+5n^2.
\]
\end{definition}

\begin{lemma}\label{lem:norm-mult}
For all $\alpha,\beta\in\Z[\sqrt{-5}]$,
\[
  \boxed{\Nm(\alpha\beta)=\Nm(\alpha)\Nm(\beta).}
\]
\end{lemma}
\begin{proof}
Write $\alpha=m+n\sqrt{-5}$ and $\beta=m'+n'\sqrt{-5}$. Then
\[
  \alpha\beta=(mm'-5nn')+(mn'+nm')\sqrt{-5},
\]
so
\[
  \Nm(\alpha\beta)=(mm'-5nn')^2+5(mn'+nm')^2.
\]
Expanding, the cross terms $\mp10mm'nn'$ cancel, and this equals
\[
  m^2m'^2+25n^2n'^2+5m^2n'^2+5n^2m'^2
  =(m^2+5n^2)\bigl(m'^2+5n'^2\bigr)=\Nm(\alpha)\Nm(\beta).
\]
\end{proof}

\begin{example}[$2$ is irreducible]
We have $\Nm(2)=\Nm(2+0\cdot\sqrt{-5})=4$. Suppose $2=\alpha\beta$.
Then
\[
  4=\Nm(2)=\Nm(\alpha)\Nm(\beta),
\]
so $\{\Nm(\alpha),\Nm(\beta)\}$ is $\{1,4\}$ or $\{2,2\}$. But
$m^2+5n^2=2$ is impossible in integers, so no element of
$\Z[\sqrt{-5}]$ has norm $2$. Hence one of $\alpha,\beta$ has norm
$1$, and is therefore $\pm1$, a unit.
\end{example}

% NOTE: the handwritten lecture ends with the setup "alpha*beta = 1
% implies N(alpha)N(beta) = 1"; the conclusion and the closing remark
% for 3 and 1 +/- sqrt(-5) complete the argument as outlined on p. 9.
Similarly, we can classify the units in $\Z[\sqrt{-5}]$. If
$\alpha\beta=1$, then
\[
  \Nm(\alpha)\Nm(\beta)=\Nm(1)=1,
\]
so $\Nm(\alpha)=m^2+5n^2=1$. This forces $n=0$ and $m=\pm1$; thus
the units of $\Z[\sqrt{-5}]$ are exactly $\pm1$.

\begin{remark}
The same argument applies to the other factors. Since
$\Nm(3)=9$ and $\Nm(1\pm\sqrt{-5})=6$, a nontrivial factorization
would require an element of norm $2$ or $3$, and $m^2+5n^2$ takes
neither value. So all four elements are irreducible. Because the
only units are $\pm1$, neither $2$ nor $3$ is an associate of
$1\pm\sqrt{-5}$. Therefore $\Z[\sqrt{-5}]$ is not a UFD.
\end{remark}

\end{document}
